\scanpage{172}
\mspage{7}
% archive sheet number: 171
connexes \unsure{sont|soient} les pts géom. déduits (via l'extension fixée $\bar K$ de $K$) de pts $x_{i\alpha}$ $(i=0,1,2)$ \deleted{des} \uline{rationnels} /$K$.
On utilise loc. cit. pour déterminer $\pi_1(X^c_{\bar K},\xi)$ comme gpe engendré par \inserted{les} $\pi_1(X^c_{0\bar K},\xi_0^c)$ et des générateurs $\unsure{g_{c'}|\gamma_{c'}}$ (indexés par l'ens des composantes connexes de $X_1$, soumis à des relations qui ne nous importent pas ici.
\deleted{Le choix \inserted{particuliers faits} des pts bases \unsure{implique}}

\deleted{Il suffira que} Cette description est canonique en dehors du choix de \deleted{\unsure{certains}} ``classes de chemins'' reliant des pts géométriques \inserted{$\bar a,\bar b$} d'une même composante connexe \inserted{$X^c_{0\bar K}$} de $X_{0\bar K}$, \deleted{\unsure{le}} [savoir le pt marqué de $X^c_{0\bar K}$ et l'image d'un pt marqué d'un $X^{c'}_{1\bar K}$ par l'une ou l'autre projection.
\deleted{\unsure{Le choix des}}
\deleted{Alors les \unsure{faits décrits ci-dessus}}
Plus précisément, on doit faire le choix d'un isom entre les foncteurs fibres $F_{\bar a},F_{\bar b}$ associés à $\bar a,\bar b$ sur la catégorie $\operatorname{Rev\,ét}(X^{\unsure{c'|c}}_{0\bar K})$ des rev. étales de $X^c_{0\bar K}=:Y_{\bar K}$.
\inserted{par construction, $\bar a,\bar b$ proviennent de pts \deleted{\unsure{géom. de $X$} $\unsure{\xi}(K)$, de} $a,b\in Y_K(K)$}
\deleted{\unsure{Au}} \unsure{Le} groupe d'inertie $I=\operatorname{Gal}(\bar K/K)$ opère par
\deleted{suite sur les foncteurs fibres $F_{\bar a},F_{\bar b}$ [donnant naissance aux opérations de $I$ sur $\pi_1(X^{c'}_{1\bar K},\bar a)$ et $\pi_1(X^{c'}_{1\bar K},\bar b)$}
\deleted{Ceci dit, on peut choisir les isom $F_{\bar a}\simeq F_{\bar b}$ de façon à commuter à l'action de \deleted{\unsure{\pi}} $I$}
sur l'ens \deleted{couples}
\[
\deleted{(\operatorname{Rev\,ét}(Y_{\bar K}),F_{\bar a}),
(\operatorname{Rev\,ét}(Y_{\bar K})F_{\bar b})}
\]
\unsure{des} ``classes de chemins'' de $\bar a$ à $\bar b$, \deleted{Nous admettrons} \inserted{qui est} \inserted{un torseur sous $\pi_1(Y_{\bar K},\bar a)$}, à groupe d'opérateurs $I$ [compatible \deleted{\unsure{avec les homomorphismes}} avec les opérations de $I$ sur les $\pi_1$).

% Cancellation scribble in the lower left margin.
\deleted{\uline{Lemme} 2.1.1. \unsure{Il existe|Soit} $Y_K$ de type fini géom connexe sur $K$, muni de $a,b\in Y_K(K)$, alors $I$ opère sur l'ens \inserted{$P$} des classes de chemins de $\bar a$ en $\bar b$ sur $Y_{\bar K}$. Ceci posé}
